\begin{table}[H]
\centering
\caption{Summary Statistics}
\label{tab:summary}
\begin{threeparttable}
\begin{tabular}{lcccc}
\toprule
& \multicolumn{2}{c}{Pre-Period (2018--2019)} & \multicolumn{2}{c}{Post-Period (2022--2023)} \\
\cmidrule(lr){2-3} \cmidrule(lr){4-5}
Variable & Mean & SD & Mean & SD \\
\midrule
Food insecure & 0.107 & 0.309 & 0.131 & 0.338 \\
Very low food security & 0.041 & 0.199 & 0.051 & 0.221 \\
SNAP receipt & 0.078 & 0.268 & 0.094 & 0.292 \\
Age & 51.836 & 17.342 & 52.150 & 17.421 \\
College degree & 0.381 & 0.486 & 0.406 & 0.491 \\
Black & 0.125 & 0.330 & 0.128 & 0.334 \\
Hispanic & 0.137 & 0.344 & 0.145 & 0.352 \\
Household size & 2.437 & 1.420 & 2.399 & 1.397 \\
Has children & 0.260 & 0.438 & 0.250 & 0.433 \\
Low income (< \$25K) & 0.195 & 0.396 & 0.158 & 0.365 \\
State SNAP rate (2019) & 0.107 & 0.022 & 0.107 & 0.022 \\
\midrule
Observations & \multicolumn{2}{c}{71,452} & \multicolumn{2}{c}{62,708} \\
\bottomrule
\end{tabular}
\begin{tablenotes}[flushleft]
\small
\item \textit{Notes:} CPS Food Security Supplement, December waves. N = 134,160 households. Weighted using CPS household weights. Pre-period: December 2018--2019 (before TFP revision and COVID). Post-period: December 2022--2023 (after TFP revision, Emergency Allotments ended). Food insecure = low or very low food security (HRFS12M1 $\geq$ 2). State SNAP rate is the 2019 ACS household SNAP participation rate, used as treatment intensity.
\end{tablenotes}
\end{threeparttable}
\end{table}
